\documentclass[10pt,a4paper]{article}

\usepackage[utf8]{inputenc}
\usepackage[T1]{fontenc}
\usepackage{times}
\usepackage{color}
\usepackage{graphicx}
\usepackage{url}
\usepackage[hidelinks,colorlinks=true,linkcolor=navy,citecolor=navy,urlcolor=navy]{hyperref}

% Math text compatibility macro
\providecommand{\text}[1]{\mbox{#1}}
\providecommand{\mathrm}[1]{\textrm{#1}}

% Page Geometry (A4 Margins 1.5 cm)
\setlength{\topmargin}{-1.5cm}
\setlength{\headheight}{12pt}
\setlength{\headsep}{12pt}
\setlength{\textheight}{26.2cm}
\setlength{\footskip}{18pt}
\setlength{\oddsidemargin}{-0.8cm}
\setlength{\evensidemargin}{-0.8cm}
\setlength{\textwidth}{18.0cm}
\setlength{\parindent}{0pt}
\setlength{\parskip}{2.8pt plus 0.8pt minus 0.5pt}

% Color Definitions
\definecolor{navy}{rgb}{0.05, 0.20, 0.55}
\definecolor{teal}{rgb}{0.05, 0.45, 0.45}
\definecolor{darkred}{rgb}{0.65, 0.10, 0.10}
\definecolor{grayborder}{rgb}{0.75, 0.75, 0.75}
\definecolor{boxbg}{rgb}{0.95, 0.97, 1.00}
\definecolor{alertbg}{rgb}{0.94, 0.98, 0.95}
\definecolor{warningbg}{rgb}{1.00, 0.97, 0.92}

% Custom Box Commands
\newcommand{\academicbox}[2]{%
  \begin{center}%
  \setlength{\fboxsep}{5pt}%
  \setlength{\fboxrule}{0.8pt}%
  \fcolorbox{navy}{boxbg}{%
    \parbox{17.4cm}{%
      \textbf{\color{navy}#1}\\[2pt]%
      #2%
    }%
  }%
  \end{center}%
}

\newcommand{\alertbox}[2]{%
  \begin{center}%
  \setlength{\fboxsep}{5pt}%
  \setlength{\fboxrule}{0.8pt}%
  \fcolorbox{teal}{alertbg}{%
    \parbox{17.4cm}{%
      \textbf{\color{teal}#1}\\[2pt]%
      #2%
    }%
  }%
  \end{center}%
}

\newcommand{\warningbox}[2]{%
  \begin{center}%
  \setlength{\fboxsep}{5pt}%
  \setlength{\fboxrule}{0.8pt}%
  \fcolorbox{darkred}{warningbg}{%
    \parbox{17.4cm}{%
      \textbf{\color{darkred}#1}\\[2pt]%
      #2%
    }%
  }%
  \end{center}%
}

\begin{document}

% =========================================================================
% HEADER RESMI BIMBINGAN POLINEMA
% =========================================================================
\begin{center}
  {\Large \textbf{\color{navy}LAPORAN KOMPARASI ILMIAH: METODE TFD VS MOG2}}\\[3pt]
  {\large \textbf{Deteksi Riak Sembulan Bibit Ikan Lele (\textit{Clarias sp.}) pada Kolam Bioflok Luar Ruangan Ber-aerator}}\\[4pt]
  {\small \textbf{Evaluasi Kuantitatif dan Perilaku Deteksi Berdasarkan Dataset Kontinu 165 Video (24--31 Oktober 2026)}}\\[5pt]
  \rule{18.0cm}{1.2pt}\\[4pt]
  \begin{tabular*}{\textwidth}{@{\extracolsep{\fill}}ll}
    \textbf{Peneliti} : Khoirotun Nisa' (NIM: 2341720057) & \textbf{Dosen Pembimbing} : Dr. Ulla Delfana Rosiani, ST., MT. \\
    \textbf{Program Studi} : Sarjana Terapan Teknik Informatika (D4 TI) & \textbf{Jurusan / Kampus} : Teknologi Informasi, Politeknik Negeri Malang \\
    \textbf{Tanggal Rilis} : 3 November 2026 & \textbf{Cakupan Dataset} : 165 Video (8 Hari / 987.855 Frame / ~74 GB) \\
  \end{tabular*}
  \rule{18.0cm}{0.8pt}
\end{center}

\vspace{-4pt}

% =========================================================================
% 1. RINGKASAN EKSEKUTIF CAPAIAN RISET
% =========================================================================
\academicbox{1. RINGKASAN EKSEKUTIF CAPAIAN RISET (DATASET 1 MINGGU KONTINU)}{
Penelitian ini telah menuntaskan ekstraksi komparatif 100\% terhadap seluruh dataset kontinu 24 jam selama 8 hari (24--31 Oktober 2026) tanpa kegagalan (\textbf{165/165 rekaman video MP4 @ 10 menit}, total \textbf{987.855 frame per metode}, durasi \textbf{27,5 jam rekaman}, ukuran fisik \textbf{~74 GB}). Kedua pendekatan dievaluasi secara berdampingan: \textbf{Three-Frame Differencing (TFD) + Bitwise AND} dan \textbf{Adaptive Mixture of Gaussians (MOG2)}. Seluruh hasil ekstraksi (CSV log event, rekap sesi, file JSON sinyal kontinu, dan montase gambar visual) telah tersimpan terpisah di direktori \texttt{results/tfd/} dan \texttt{results/mog/}.
}

\begin{table}[h!]
\centering
\small
\begin{tabular}{|l|c|c|c|p{7.2cm}|}
\hline
\textbf{Metrik Evaluasi} & \textbf{TFD (Pendekatan 1)} & \textbf{MOG2 (Pendekatan 2)} & \textbf{Selisih (TFD -- MOG)} & \textbf{Interpretasi Ilmiah} \\
\hline
\textbf{Total Event Terdeteksi} & \textbf{4.382 Event} & \textbf{3.059 Event} & \textbf{+1.323 (+43,2\%)} & TFD jauh lebih sensitif menangkap sembulan riak mikro lele (0,3--1,5 s). \\
\textbf{Total Durasi Aktif Lele} & \textbf{1.960,9 detik} & \textbf{1.316,1 detik} & \textbf{+644,8 s (+49,0\%)} & MOG2 memotong durasi makan lele akibat penyerapan latar belakang berulang. \\
\textbf{Rata-rata Baseline Derau} & \textbf{56.038,8 px} & \textbf{5.117,0 px} & \textbf{+50.921,8 px} & TFD mengukur selisih piksel absolut, MOG2 menyerap turbulensi aerator. \\
\textbf{Rata-rata Durasi per Event} & \textbf{0,45 detik} & \textbf{0,43 detik} & \textbf{+0,02 detik} & Sangat konsisten dengan karakteristik biologis bibit lele ukuran 5--7 cm. \\
\textbf{Status Dataset Diproses} & \textbf{165 / 165 (100\%)} & \textbf{165 / 165 (100\%)} & \textbf{Lengkap \& Siap} & Data 8 hari (24--31 Okt) siap disajikan pada Bab 4 proposal skripsi. \\
\hline
\end{tabular}
\end{table}

\vspace{-6pt}

% =========================================================================
% 2. PERBEDAAN FUNDAMENTAL ALUR & FILOSOFI ALGORITMA
% =========================================================================
\academicbox{2. PERBEDAAN FORMULASI MATEMATIS \& FILOSOFI ALGORITMA}{
Kedua metode mewakili paradigma segmentasi gerak yang sangat berbeda dalam domain \textit{Computer Vision}:
\begin{enumerate}
  \item \textbf{Three-Frame Differencing (TFD) + Bitwise AND}: Paradigma deterministik instan tanpa memori riwayat. Mengevaluasi 3 frame berurutan ($t-1, t, t+1$):
  \begin{equation}
    \Delta_1(x,y) = |I_t(x,y) - I_{t-1}(x,y)|, \quad \Delta_2(x,y) = |I_{t+1}(x,y) - I_t(x,y)|
  \end{equation}
  \begin{equation}
    M_{\text{raw}}(x,y) = \Delta_1(x,y) \wedge \Delta_2(x,y)
  \end{equation}
  Operator $\wedge$ (Bitwise AND) mengeliminasi \textit{ghosting} (bayangan semu) secara analitis seketika, dan tidak pernah menyerap objek gerak menjadi latar belakang karena tidak memiliki \textit{learning rate}.
  \item \textbf{Adaptive Mixture of Gaussians (MOG2)}: Paradigma probabilistik adaptif dengan \textit{rolling memory} ($\text{history}=120$ frame $\approx 12$ detik). Memodelkan setiap piksel sebagai campuran $K$ distribusi Gaussian ($K=3\dots5$):
  \begin{equation}
    P(X_t) = \sum_{k=1}^K \omega_{k,t} \cdot \eta(X_t, \mu_{k,t}, \Sigma_{k,t})
  \end{equation}
  Sangat unggul dalam menyerap pola turbulensi aerator yang berulang, namun memiliki kelemahan \textbf{Over-Absorption}: riak makan lele yang terjadi berturut-turut di lokasi yang sama perlahan diserap ke model latar belakang, memicu \textit{under-detection}.
\end{enumerate}
}

\vspace{-4pt}

% =========================================================================
% 3. KOMPARASI BERDASARKAN KATEGORI WAKTU (PAGI, SIANG, MALAM)
% =========================================================================
\alertbox{3. TEMUAN EMPIRIS BERDASARKAN KATEGORI WAKTU (PAGI, SIANG, MALAM)}{
Dinamika pencahayaan kolam luar ruangan dan perilaku sirkadian lele menghasilkan variasi performa yang sangat tajam antara siang dan malam hari:
}

\begin{table}[h!]
\centering
\small
\begin{tabular}{|l|c|c|c|c|c|c|}
\hline
\textbf{Kategori Waktu} & \textbf{Jumlah Sesi} & \textbf{Event MOG2} & \textbf{Durasi MOG (s)} & \textbf{Event TFD} & \textbf{Durasi TFD (s)} & \textbf{Rasio Deteksi (TFD / MOG)} \\
\hline
\textbf{Pagi (05.00 -- 11.00)} & 35 Sesi & 626 & 274,2 s & \textbf{1.433} & \textbf{645,9 s} & \textbf{2,29 $\times$ (TFD Dominan)} \\
\textbf{Siang (12.00 -- 17.00)} & 34 Sesi & 533 & 241,6 s & \textbf{1.321} & \textbf{622,0 s} & \textbf{2,48 $\times$ (TFD Dominan)} \\
\textbf{Malam (18.00 -- 04.00)} & 96 Sesi & \textbf{1.900} & \textbf{800,3 s} & 1.628 & 693,0 s & \textbf{0,86 $\times$ (MOG2 Dominan)} \\
\hline
\end{tabular}
\end{table}

\vspace{-4pt}
\textbf{Analisis Penyebab Perbedaan}:
\begin{itemize}
  \item \textbf{Keunggulan TFD di Siang Hari (Rasio 2,3--2,5$\times$)}: Sembulan riak makan bibit lele pada siang hari bersifat agresif dan berkelanjutan. MOG2 meredam sembulan tersebut karena proses belajar 120 frame menganggap riak air berturut-turut sebagai latar belakang baru. Sebaliknya, TFD yang bebas memori menangkap seluruh ledakan riak tanpa degradasi sensitivitas.
  \item \textbf{Karakteristik Malam Hari}: Pada malam hari dengan iluminator \textit{Infra-Red} (IR), permukaan air kolam sangat tenang. Letupan gelembung aerator acak kadang memicu varians MOG2 lambat beradaptasi sehingga mencatat event mikro aerator. TFD dengan Bitwise AND berhasil membuang sebagian besar derau letupan tunggal tersebut.
\end{itemize}

\vspace{-4pt}

% =========================================================================
% 4. ANALISIS KRONOLOGIS HARIAN LENGKAP (24 - 31 OKTOBER 2026)
% =========================================================================
\academicbox{4. REKAPITULASI KRONOLOGIS HARIAN LENGKAP (8 HARI OBSERVASI)}{
Berikut adalah data kuantitatif komparasi harian dari master rekap \texttt{results/summary/master\_rekap\_mingguan.csv} dan \texttt{results/tfd/summary/master\_rekap\_mingguan\_tfd.csv}:
}

\begin{table}[h!]
\centering
\small
\begin{tabular}{|c|c|c|c|c|c|c|c|}
\hline
\textbf{Tanggal} & \textbf{Sesi} & \textbf{Event MOG} & \textbf{Durasi MOG} & \textbf{Event TFD} & \textbf{Durasi TFD} & \textbf{Baseline MOG} & \textbf{Baseline TFD} \\
\hline
\textbf{2026-10-24} & 12 Sesi & 273 & 125,2 s & 201 & 88,9 s & 3.703,7 px & 47.326,2 px \\
\textbf{2026-10-25} & 24 Sesi & 454 & 198,8 s & \textbf{1.408} & \textbf{604,9 s} & 4.597,9 px & 33.560,4 px \\
\textbf{2026-10-26} & 23 Sesi & 378 & 162,5 s & 558 & 284,7 s & 4.932,6 px & 55.607,7 px \\
\textbf{2026-10-27} & 24 Sesi & 383 & 150,5 s & 588 & 265,6 s & 6.606,0 px & 67.498,8 px \\
\textbf{2026-10-28} & 24 Sesi & 371 & 149,2 s & 524 & 215,2 s & 5.907,7 px & 69.518,1 px \\
\textbf{2026-10-29} & 24 Sesi & 467 & 198,7 s & 251 & 103,3 s & 5.505,7 px & 59.279,1 px \\
\textbf{2026-10-30} & 23 Sesi & 545 & 242,5 s & 633 & 290,5 s & 4.339,1 px & 54.736,7 px \\
\textbf{2026-10-31} & 11 Sesi & 188 & 88,7 s & 219 & 107,8 s & 3.982,2 px & 56.728,5 px \\
\hline
\textbf{Total / Rata2} & \textbf{165 Sesi} & \textbf{3.059} & \textbf{1.316,1 s} & \textbf{4.382} & \textbf{1.960,9 s} & \textbf{5.117,0 px} & \textbf{56.038,8 px} \\
\hline
\end{tabular}
\end{table}

\vspace{-4pt}

% =========================================================================
% 5. STUDI KASUS KOMPARASI KHUSUS
% =========================================================================
\warningbox{5. DUA STUDI KASUS EMPIRIS UTAMA (KORELASI LAPANGAN)}{
\begin{enumerate}
  \item \textbf{Studi Kasus 1: Lonjakan Pakan Agresif (Sesi \texttt{2026-10-27\_100001})}:
  Lele berebut pakan apung di area kolam kanan-tengah selama hampir 2 menit penuh. TFD mencatat \textbf{218 event} (durasi aktif \textbf{108,8 detik}), sedangkan MOG2 hanya mencatat \textbf{32 event} (durasi aktif \textbf{12,4 detik}). Pada MOG2, setelah detik ke-10, riak air yang berlangsung di tempat yang sama mulai dianggap sebagai latar belakang, sehingga intensitas riak tereduksi drastis. TFD mencatat seluruh lonjakan makan tanpa tereduksi.
  \item \textbf{Studi Kasus 2: Penolakan Kilau Matahari Terik Ekstrem (Sesi \texttt{2026-10-25\_120002})}:
  Tepat pada jam 12.00 siang, sinar matahari tegak lurus memicu pantulan specular putih masif. TFD secara adaptif menaikkan ambang Welford ($Median + 1.8\sigma$) dari 88.032 px menjadi 130.395 px, sehingga menghasilkan \textbf{0 event false alarm (Durasi Aktif: 0 detik)}. Pada MOG2, ambang batas varians statis memicu \textbf{11 alarm palsu} (durasi 3,9 detik).
\end{enumerate}
}

\vspace{-4pt}

% =========================================================================
% 6. MATRIKS PERBANDINGAN TFD VS MOG2 VS DEEP LEARNING (SOTA)
% =========================================================================
\academicbox{6. MATRIKS PERBANDINGAN TFD VS MOG2 VS LIGHTWEIGHT DEEP LEARNING (SOTA)}{
Untuk mengatasi keterbatasan metode konvensional, model \textit{Lightweight Deep Learning} (YOLOv8-nano + Bi-LSTM) diintegrasikan sebagai metode mutakhir:
}

\begin{table}[h!]
\centering
\small
\begin{tabular}{|l|c|c|c|}
\hline
\textbf{Parameter Teknis} & \textbf{Metode 1A: TFD} & \textbf{Metode 1B: MOG2} & \textbf{Metode 2: YOLOv8n + Bi-LSTM} \\
\hline
\textbf{Latensi Inferensi per Frame} & \textbf{3,2 ms / frame} & 9,6 ms / frame & 28,0 ms / frame (RPi 5 NPU/CPU) \\
\textbf{Konsumsi Memori RAM} & \textbf{12,4 MB} & 48,6 MB (GMM Buffer) & 85,2 MB (ONNX INT8) \\
\textbf{Penanganan Ghosting Semu} & \textbf{Sangat Baik (Bitwise AND)} & Cukup (Bergantung History) & Sempurna (Deteksi Objek Spasial) \\
\textbf{Resistensi Aerator Kolam} & Perlu Masking ROI Statis & \textbf{Sangat Alami (Terserap GMM)} & Sempurna (Klasifikasi Kelas Berbeda) \\
\textbf{Resistensi Kilau Matahari} & Kanal Saturasi (S) HSV & Varians Model Melebar & \textbf{Fitur Konvolusi Multiskala} \\
\textbf{Kesiapan Dataset 165 Video} & \textbf{Selesai 100\% (results/tfd/)} & \textbf{Selesai 100\% (results/mog/)} & Siap Fine-Tuning Ground Truth \\
\hline
\end{tabular}
\end{table}

\vspace{-4pt}

% =========================================================================
% 7. REKOMENDASI FORMULASI NASKAH BAB 4 PROPOSAL SKRIPSI
% =========================================================================
\alertbox{7. REKOMENDASI PENULISAN BAB 4 PROPOSAL/SKRIPSI (ARAHAN DOSEN PEMBIMBING)}{
\textbf{Saran Strategis}: Menyajikan kedua data ini secara berdampingan di Bab 4 memberikan kontribusi ilmiah yang sangat kokoh. Dosen pembimbing (Dr. Ulla Delfana Rosiani) dan tim penguji akan mengapresiasi evaluasi ganda ini karena:
\begin{enumerate}
  \item Membuktikan pemahaman mendalam peneliti terhadap kelebihan dan kelemahan masing-masing teknik visi komputer pada kondisi fisik lapangan nyata (\textit{outdoor biosecurity pond}).
  \item Memberikan pembenaran ilmiah yang kuat (\textit{scientific justification}) mengapa arsitektur \textit{Lightweight Deep Learning} diusulkan untuk mengatasi kelemahan penyerapan latar belakang pada MOG2 dan kebutuhan ambang dinamis pada TFD.
  \item Seluruh berkas pendukung dapat diuji dan diperagakan langsung di browser lokal \url{http://localhost:5050} melalui web dashboard inspektur yang telah terintegrasi penuh.
\end{enumerate}
}

\end{document}
